\section{Introduction}\label{sec:introduction}

Large corpora are indispensable in contemporary political science, but the objects we ultimately analyze are not the raw texts themselves. What researchers actually want are structured signals: whether an article reports a protest event, which actor sponsored an amendment, the date and location of a strike, or the ideological position of a party manifesto on decentralization or environmental policy. Text is rarely the endpoint. It is the raw material from which we hope to extract something analytically useful, and the central methodological question is how to move from messy, unstructured documents to the clean, interpretable variables that enter our regressions and substantive analyses.

This challenge has driven decades of innovation, from hand-coding schemes to dictionary methods to topic models. More recently, large language models have emerged as a powerful tool for structured extraction, and their use is now ubiquitous across the social sciences. A striking example appears in recent work by \citet{BenoitEtAl2025}, who use ensembles of LLMs to estimate party positions on key policy dimensions from manifestos written in dozens of languages. Their method asks LLMs to first summarize what a text says about particular issues, then to score the author's positions on those issues. The results correlate highly with expert surveys and, when applied to coalition policy declarations, align more closely with standard models of government formation than traditional hand-coded estimates. This is ``natural language understanding'' in the applied sense: using machines to replicate the outputs of human interpretive judgments, at scale, cheaply, and in practically any language.

Why the excitement? Qualitative methods that rely on human experts are costly and do not scale. Quantitative text-as-data approaches scale well but often rest on strong assumptions about what word frequencies or topic proportions actually capture. LLM-based methods promise something in between: the depth of traditional qualitative analysis combined with the scalability of computational approaches. But this promise comes with a methodological caveat. Even when LLM predictions are highly accurate in a classification sense, using them directly as input variables in downstream analyses can introduce bias and invalidate confidence intervals in ways that are not immediately apparent.

Recent work on Design-based Supervised Learning \citep{EgamiEtAl2024} makes this concrete. Even when prediction errors are small (say, 5--10\% misclassification), naively substituting LLM-generated labels for ground-truth annotations can substantially bias downstream regression coefficients and produce confidence intervals with far less than nominal coverage. The problem is not that the LLM is ``wrong'' in some average sense, but that its errors may be systematic in ways that correlate with the very outcomes we are trying to explain. The DSL framework addresses this by combining surrogate labels with a smaller number of gold-standard human annotations, using a doubly-robust procedure that guarantees valid inference even when the surrogate is arbitrarily biased. If we want to use LLM outputs as data, we need formal methods for understanding and correcting for their relationship to the ground truth.

What both of these literatures point toward, albeit from different angles, is the idea of an underlying \emph{oracle} judgment. In \citet{BenoitEtAl2025}, the oracle is implicit: the expert survey scores that serve as the benchmark against which LLM estimates are validated. In \citet{EgamiEtAl2024}, it is explicit: the gold-standard labels that anchor the debiasing procedure. In both cases, the fundamental object of interest is not the LLM output per se, but some latent ``true'' value that the researcher actually cares about and that expensive human judgment can, in principle, recover.

\subsection{The Oracle-Preserving Summarization Framework}

We formalize this intuition by positing an oracle $\fstar$ that maps a full document $x$ to the task value $\fstar(x) \in \Y$ that the researcher actually cares about. In this view, LLMs are useful only insofar as they preserve (or at least accurately predict) the information contained in $\fstar(x)$. Summaries, scores, and extracted variables are valid research objects only if they leave the oracle untouched. We therefore treat the summarizer $g$ as a targeted compression operator: its sole job is to discard tokens that are provably irrelevant to $\fstar$ while preserving every bit of evidence needed for the oracle to reach the same verdict.\footnote{Put differently, $g(x)$ must be a sufficient statistic for $\fstar(x)$. It may be terse or awkward, but it cannot omit the features that $\fstar$ inspects.} The oracle might be human annotation, a high-quality-but-expensive model, or a hypothetical infinite-context extractor.\footnote{In practice we approximate $\fstar$ with whatever expensive-but-accessible supervision we can afford: expert coders, deliberation-grade models, or hybrid workflows. The audit is therefore relative to that chosen oracle.} It is often too costly or infeasible to apply $\fstar$ to every raw document, which is precisely why we rely on compressed representations $g(x)$ in the first place.

This framing provides a natural bridge to two broader literatures. First, we are asking for the summary to be a \emph{sufficient statistic} for $\fstar(x)$ in the classical sense from statistical decision theory \citep{Blackwell1953}: once you have the summary, the original document provides no additional information about the task value. Any downstream analysis that depends only on $\fstar$ will give the same answer whether we use the summary or the full document. Second, we connect to work on \emph{mergeable summaries} in database theory \citep{agarwal2013mergeable}, which asks when compact sketches of data can be combined without losing information about some target statistic. In the mergeable-summaries literature, one constructs explicit algebraic data structures (sketches) that preserve commutative statistics like counts or heavy hitters under merge operations. The key operation is a hand-designed merge function that combines two sketches while maintaining the invariant. Our setting differs in two ways: we work with string concatenation, which is non-commutative (order matters), and we seek to preserve arbitrary semantic oracles like event extraction or narrative flow rather than simple numeric aggregates. We therefore replace the explicit algebra with \emph{learned} local consistency conditions that force a stochastic summarizer to behave \emph{as if} it were a mergeable sketch. When the oracle happens to be symmetric (e.g., counting events where order is irrelevant), our conditions collapse exactly to standard mergeability; Appendix~\ref{app:mergeable} makes this correspondence precise.

\subsection{Three Contributions}

This paper makes three contributions, each aimed at a different audience but together forming a unified framework for reliable semantic extraction at scale.

\paragraph{Contribution 1: Theory.} We introduce three consistency conditions---Sufficiency (C1), Idempotence (C2), and Merge Consistency (C3)---that can be checked locally at each node of a summary tree. When these conditions hold, the final summary preserves the oracle in expectation: $\E[d_\Y(\fstar(Z), \fstar(x))] = 0$, where $Z$ is the summary and $d_\Y$ is a metric on the oracle space. Crucially, these checks require only the immediate inputs at each node, not the full document. We prove schedule invariance (any binary tree over the same blocks yields the same expected oracle), fold-of-folds invariance (two-level hierarchies are interchangeable), and multi-round preservation (additional cleanup passes maintain the guarantee). We also prove a tree-level optimizer-transfer theorem for stochastic adaptive objectives: when the oracle itself may be imperfect, optimization error decomposes into expected oracle uncertainty plus expected tree-transport error. All proofs are \emph{machine-verified} in Lean~4, comprising more than 50,000 lines of code across 123 files. The formalization requires only a single axiom---Expected Lipschitz continuity of group losses under Random Utility Models---which we justify via McFadden's discrete choice theory \citep{McFadden1974}.

\paragraph{Contribution 2: Software.} We release \textsc{ThinkingTrees}, an open-source Python package implementing the full OPS pipeline. The package provides:
\begin{itemize}[nosep]
    \item \textbf{TreeBuilder}: Async-first bottom-up tree construction with configurable chunking
    \item \textbf{TreeAuditor}: Probabilistic verification using Hoeffding bounds
    \item \textbf{DSPy Integration}: Optimizers (GEPA, MIPRO, Bootstrap) for prompt tuning
    \item \textbf{Preference Learning}: DPO and GRPO training on oracle-preserving summaries
    \item \textbf{Plugin Architecture}: Extensible tasks (Manifesto RILE, Congressional Record) and datasets
\end{itemize}
The codebase comprises 182,000+ lines of Python with comprehensive test coverage. Section~\ref{sec:software} describes the architecture; Appendix~\ref{app:software} provides API documentation.

\paragraph{Contribution 3: Empirics.} We demonstrate OPS on two applications central to political science:
\begin{enumerate}[nosep]
    \item \textbf{Manifesto Project RILE Scoring}: Estimating party positions on a left-right scale ($-100$ to $+100$) from party manifestos, comparing OPS-compressed summaries to direct full-document scoring (Section~\ref{sec:manifesto}).
    \item \textbf{Congressional Record Policy Positions}: Extracting issue-specific positions (immigration, climate, healthcare) from long Congressional speeches, demonstrating OPS on multi-dimensional coding tasks (Section~\ref{sec:congressional}).
\end{enumerate}
Both applications use explicit rubrics that define what information the oracle must preserve, enabling rigorous auditing. Appendix~\ref{app:empirical} provides full results and robustness checks.

\subsection{Why Hierarchical Summarization?}

Two practical constraints motivate the hierarchical approach. First, many documents are longer than a model's context window. In practice, one partitions $x$ into contiguous blocks, summarizes each block with an LLM, and recursively merges those summaries up a tree. Second, even when context is not the binding constraint, researchers often \emph{choose} to summarize for cost, latency, or workflow reasons: storing compact ``event summaries'' that can be relabeled as coding rules evolve, or generating intermediate representations that multiple downstream tasks can query. Either way, hierarchical summarization changes the object of analysis, and a natural question arises: \emph{How can we ensure this pipeline leaves $\fstar(x)$ unchanged on average, without running the expensive oracle on the full text?}

We answer this question by introducing three consistency conditions (leaf-level sufficiency, on-range idempotence, and merge consistency) that can be checked locally, at each node of the summary tree, without ever processing the full document. When these conditions hold, the final summary preserves the oracle in expectation. When the conditions are violated, we can estimate how often and by how much. Crucially, these same conditions double as optimization objectives. Using DSPy \citep{khattab2024dspy}, we treat the summarization policy as a programmable system whose prompt parameters are updated until violations disappear. Because each check inspects only the immediate child summaries, we obtain cheap supervision signals without requiring new external labels. The audits that certify a pipeline also describe how to tune it.

\subsection{Paper Roadmap}

The remainder of the paper develops this framework across five parts:

\begin{description}[nosep]
    \item[Part I: Core Framework] Section~\ref{sec:framework} introduces notation and the three consistency conditions. Section~\ref{sec:theorems} states the main theorems (with proofs in Appendix~\ref{app:proofs-core}). Section~\ref{sec:audit} develops the probabilistic audit.

    \item[Part II: Software] Section~\ref{sec:software} describes the \textsc{ThinkingTrees} package. Section~\ref{sec:preference} connects OPS to preference learning (DPO, GRPO).

    \item[Part III: Applications] Section~\ref{sec:manifesto} applies OPS to Manifesto Project RILE scoring. Section~\ref{sec:congressional} applies it to Congressional Record policy extraction.

    \item[Part IV: Verification] Section~\ref{sec:verification} documents the Lean~4 formalization.

    \item[Part V: Discussion] Section~\ref{sec:related} surveys related work. Section~\ref{sec:conclusion} concludes.
\end{description}

Throughout, we focus on decomposable, information-centric tasks: event extraction, coding tuples, structured aggregation, and guided abstractive summaries with explicit rubrics where the oracle space $\Y$ is well-defined on bounded spans of text. Open-ended creative summarization, where the oracle rewards prose quality directly, falls outside the scope of this paper. So too do settings where the LLM itself introduces systematic bias unrelated to the summarization process. If a model were politically slanted or subject to content filtering that distorts the underlying signal, even a perfect summarization pipeline would faithfully preserve a corrupted oracle, and additional diagnostics would be needed.
